\section{The first stable level can be prescribed}

Fix a prime $p\equiv3\pmod4$ and an integer $k\ge1$. On
$(\mathbb Z/p^n\mathbb Z)^2$ take one step in each of the directions
$p^k e_1$ and $e_2$, and $p$ copies of the step $e_1$. Each step from a
vertex contributes one undirected edge. Parallel edges count separately
and loops contribute zero to the Laplacian. These connected towers
all have the same number $p+2$ of base loops for fixed $p$.
Let $T_n$ be the $p$-adic valuation of the spanning-tree count.

\begin{theorem}
Put $s=\min(k,n)$ and $C=2(p+1)/(p-1)$. At every level $n\ge1$,
\[
 T_n=2s p^n+C(p^n-p^s)+(p^s-1)+2s-2n .
\]
The eventual polynomial is
\[
 P_k(X,Y)=(2k+C)X-2Y+(1-C)p^k+2k-1.
\]
It gives $T_n=P_k(p^n,n)$ from level $k$ onward. For $k\ge2$ it fails
at level $k-1$, so $k$ is the exact first stable level.
\end{theorem}

\begin{proof}
Write $A(U)=2-U-U^{-1}=-(U-1)^2/U$.
For a pair $(U,V)$ of $p$-power roots of unity of joint conductor $p^j>1$,
let $e_j=(p-1)p^{j-1}$. In a common primitive coordinate $\zeta$ of order
$p^j$, write $U=\zeta^a$, $V=\zeta^b$, with $a,b$ not both divisible
by $p$. The reduction of
$(A(U)+A(V))/(\zeta-1)^2$ modulo $\zeta-1$ is $-(a^2+b^2)$. Since
$-1$ is a nonsquare modulo $p$, this is nonzero. Consequently
\[
 v_p(A(U)+A(V))=2/e_j\le1 .
\]
The eigenvalue at a character $(X,Y)$ is
$A(X^{p^k})+A(Y)+pA(X)$. The final term has valuation strictly above one,
with the convention $v_p(0)=+\infty$. Thus whenever
$(X^{p^k},Y)\ne(1,1)$, the joint conductor of this image alone determines
the valuation.

The map $(X,Y)\mapsto(X^{p^k},Y)$ has fibers of size $p^s$ and image
isomorphic to $C_{p^{n-s}}\times C_{p^n}$, where $C_m$ denotes a cyclic
group of order $m$. For $1\le j\le n-s$ its exact conductor-$p^j$ shell
has $p^{2j}-p^{2j-2}$ elements. Their aggregate valuation is therefore
$2(p+1)p^{j-1}$. For $n-s<j\le n$ there are $p^{n-s}e_j$ elements in
the shell, contributing $2p^{n-s}$. Summing these shell contributions
and multiplying by $p^s$ gives $C(p^n-p^s)+2s p^n$.
On the kernel with the trivial character removed, the eigenvalue is
$pA(X)$ for each nontrivial $p^s$-th root $X$. There are $p^s-1$ such
roots, and
$\prod_{X\ne1}(1-X)=p^s$ gives $\sum_{X\ne1}v_p(A(X))=2s$.
The matrix-tree denominator contributes $-2n$. This proves the formula.

For $n\le k$ it reads $T_n=(2n+1)p^n-1$, and for $n\ge k$ it equals
$P_k(p^n,n)$. In particular the discrepancy one level before $k$ is
\[
 T_{k-1}-P_k(p^{k-1},k-1)=(p+1)p^{k-1}-2>0 .
\]
Distinct polynomials in $p^n,n$ cannot agree on an infinite tail:
the term of highest degree in $p^n$ dominates all lower exponential
orders, even with polynomial factors in $n$. This establishes the
claimed onset for the unique eventual polynomial.
\end{proof}

The growing onset is visible directly in the formal specializations.
For primitive $(a,b)\in\mathbb Z_p^2$, substitute $X=(1+z)^a$ and
$Y=(1+z)^b$ in the Laplacian polynomial and reduce modulo $p$.
The first nonzero degree is
\[
 2p^{\min(k+v_p(a),v_p(b))}.
\]
Indeed the two summands that survive reduction have these respective
initial orders, and in an equal-order tie the same anisotropic sum of
squares prevents cancellation. The maximum is exactly $2p^k$, attained
at $(1,0)$. Thus even within this family with no identically zero
primitive specialization, neither the maximal degree required to see a
unit nor the onset admits a bound in terms of $p$, dimension and base
edge count alone. Replacing $p^k e_1$ by the zero step gives the law
$(2n+1)p^n-1$ at every level, despite agreement of the voltage data
modulo $p^k$.

For $p=3$ and $k=3$ the stable law is $10\cdot3^n-2n-76$ and the first
five values are $8,44,188,726,2344$. The script
\texttt{verify\_delayed\_tower.py} checks the eigenvalues by exact
cyclotomic reduction. It independently constructs integer Laplacian
cofactors for its examples with at most 81 vertices.

\paragraph{A version with distinct steps.}
For $p\ge7$ in the same congruence class, replace the $p$ copies of
$e_1$ by $(1+i p^k)e_1$ for $0\le i<p$. Every valuation stays the same.
Indeed, writing $Z=X^{p^k}$, the replacement contributes
$2p-X\Phi_p(Z)-X^{-1}\Phi_p(Z^{-1})$. At $Z=1$ this is $pA(X)$.
If $Z$ has order $p$ the two cyclotomic factors vanish. If its order is
larger, each has valuation $(p-1)v_p(Z-1)$ by
$\Phi_p(Z)=(Z^p-1)/(Z-1)$. Hence for $Z\ne1$ the new contribution has
valuation strictly greater than $2v_p(Z-1)$, which in turn is at least
$v_p(A(Z)+A(Y))$. For $n>k$ all steps are nonzero and distinct up to
sign, so the graphs at those levels are simple.
